\section{Banco de questões — Gestão e Governança de TI}

\subsection*{N1 — Fundamento competitivo}

\begin{questao}{\nivel{N1} 14.01 — Cebraspe/Certo ou Errado}
No COBIT 2019, EDM reúne objetivos de governança, enquanto APO, BAI, DSS e MEA reúnem objetivos de gestão.
\end{questao}

\begin{questao}{\nivel{N1} 14.02 — Cebraspe/Certo ou Errado}
A distinção entre governança e gestão pode ser feita apenas pelo verbo da atividade: toda atividade de ``monitorar'' pertence à governança.
\end{questao}

\begin{questao}{\nivel{N1} 14.03 — Cebraspe/Certo ou Errado}
Os fatores de desenho do COBIT 2019 indicam que o sistema de governança deve ser adaptado ao contexto da organização, e não copiado integralmente de outra entidade.
\end{questao}

\begin{questao}{\nivel{N1} 14.04 — Cebraspe/Certo ou Errado}
No ITIL 4, incident management e problem management perseguem exatamente o mesmo objetivo, diferindo apenas na nomenclatura do registro.
\end{questao}

\begin{questao}{\nivel{N1} 14.05 — Cebraspe/Certo ou Errado}
O princípio ``Start where you are'' não significa conservar tudo o que existe, mas avaliar o estado atual e aproveitar o que funciona antes de reconstruir sem evidência.
\end{questao}

\begin{questao}{\nivel{N1} 14.06 — Cebraspe/Certo ou Errado}
Apetite a risco é decisão de governança conectada aos objetivos institucionais; não deve ser definido isoladamente pela área técnica de TI.
\end{questao}

\begin{questao}{\nivel{N1} 14.07 — Cebraspe/Certo ou Errado}
Um KPI sem fórmula, fonte ou responsável pode ser visualmente útil em dashboard, mas possui baixa auditabilidade e reprodutibilidade.
\end{questao}

\begin{questao}{\nivel{N1} 14.08 — Cebraspe/Certo ou Errado}
Em BPMN, um fluxo de sequência pode atravessar pools para representar troca de mensagens entre participantes distintos.
\end{questao}

\begin{questao}{\nivel{N1} 14.09 — Cebraspe/Certo ou Errado}
Conformidade com uma norma ou certificação não implica risco zero, pois escopo, implementação e risco residual continuam relevantes.
\end{questao}

\begin{questao}{\nivel{N1} 14.10 — Cebraspe/Certo ou Errado}
Governança de dados exige definição de autoridade e responsabilidade sobre significado, qualidade e acesso aos dados, não apenas escolha de ferramentas de armazenamento.
\end{questao}

\subsection*{N2 — Aplicação}

\begin{questao}{\nivel{N2} 14.11 — FGV/COBIT}
O conselho avalia três alternativas de sourcing para serviços críticos, define critérios de valor e risco, aprova a direção estratégica e acompanha os resultados. A atividade está mais associada a:
\begin{enumerate}[label=\Alph*)]
\item EDM;
\item DSS;
\item BAI;
\item service request management;
\item operação de data center.
\end{enumerate}
\end{questao}

\begin{questao}{\nivel{N2} 14.12 — FCC/COBIT}
Uma equipe planeja arquitetura, orçamento e capacidades necessárias para cumprir objetivos corporativos de TI. O domínio COBIT mais aderente é:
\begin{enumerate}[label=\Alph*)]
\item EDM;
\item APO;
\item DSS;
\item MEA exclusivamente;
\item Incident Management.
\end{enumerate}
\end{questao}

\begin{questao}{\nivel{N2} 14.13 — Cebraspe/Certo ou Errado}
A cascata de objetivos do COBIT ajuda a conectar necessidades de stakeholders a objetivos empresariais, de alinhamento e de governança/gestão.
\end{questao}

\begin{questao}{\nivel{N2} 14.14 — FGV/ITIL}
Um serviço foi restaurado em 20 minutos, mas o mesmo incidente reaparece semanalmente. A prática que deve investigar causas e erros conhecidos é:
\begin{enumerate}[label=\Alph*)]
\item Incident management;
\item Problem management;
\item Service request management;
\item Deployment management;
\item Relationship management.
\end{enumerate}
\end{questao}

\begin{questao}{\nivel{N2} 14.15 — Cebraspe/Certo ou Errado}
Change enablement busca aumentar a taxa de mudanças bem-sucedidas por avaliação de risco e autorização apropriada; isso não exige que toda mudança passe por um mesmo comitê manual.
\end{questao}

\begin{questao}{\nivel{N2} 14.16 — FCC/ISO 31000}
No processo de gestão de riscos, identificação, análise e avaliação compõem, em conjunto:
\begin{enumerate}[label=\Alph*)]
\item tratamento de riscos;
\item avaliação de riscos;
\item comunicação apenas;
\item auditoria operacional;
\item gestão de incidentes.
\end{enumerate}
\end{questao}

\begin{questao}{\nivel{N2} 14.17 — FGV/PDTI}
Um PDTI contém 65 iniciativas, mas não informa objetivo institucional, responsável, indicador ou risco. A principal deficiência é:
\begin{enumerate}[label=\Alph*)]
\item excesso de alinhamento estratégico;
\item incapacidade de transformar estratégia em execução monitorável;
\item ausência de metodologia ágil;
\item falta de BPMN;
\item uso excessivo de KPI.
\end{enumerate}
\end{questao}

\begin{questao}{\nivel{N2} 14.18 — Cebraspe/Certo ou Errado}
Um indicador ``número de tickets fechados'' pode induzir fechamento prematuro se usado isoladamente como meta, ilustrando risco de gaming de métricas.
\end{questao}

\begin{questao}{\nivel{N2} 14.19 — FCC/BSC}
Em uma organização pública, o BSC deve ser entendido principalmente como:
\begin{enumerate}[label=\Alph*)]
\item painel exclusivamente financeiro;
\item mecanismo para traduzir estratégia em objetivos e indicadores relacionados, adaptado ao valor público;
\item substituto da gestão de riscos;
\item catálogo de serviços;
\item mapa de processos BPMN.
\end{enumerate}
\end{questao}

\begin{questao}{\nivel{N2} 14.20 — FGV/BPMN}
Dois participantes distintos trocam uma solicitação e uma resposta. Em BPMN, a comunicação entre pools deve ser representada por:
\begin{enumerate}[label=\Alph*)]
\item sequence flow;
\item message flow;
\item gateway paralelo;
\item associação de dados apenas;
\item evento de erro obrigatoriamente.
\end{enumerate}
\end{questao}

\begin{questao}{\nivel{N2} 14.21 — Cebraspe/Certo ou Errado}
Automatizar um processo ineficiente sem revisar etapas e controles pode apenas acelerar desperdício, o que é incompatível com a lógica ``Optimize and automate'' do ITIL 4.
\end{questao}

\begin{questao}{\nivel{N2} 14.22 — FCC/PMBOK 8}
Um projeto experimental de IA de duas semanas recebe a mesma documentação e governança de uma implantação nacional crítica. O princípio mais diretamente negligenciado é:
\begin{enumerate}[label=\Alph*)]
\item tailoring/adaptação ao contexto;
\item disponibilidade;
\item gestão de incidentes;
\item normalização;
\item segregação de rede.
\end{enumerate}
\end{questao}

\begin{questao}{\nivel{N2} 14.23 — Cebraspe/Certo ou Errado}
No PMBOK 8, prazo e custo continuam relevantes, mas sucesso de projeto não precisa ser reduzido à tríplice restrição se valor, benefícios e adoção forem parte do objetivo.
\end{questao}

\begin{questao}{\nivel{N2} 14.24 — FGV/LAI-LGPD}
Uma unidade pretende publicar base nominal com dados pessoais sob argumento de transparência ativa. A decisão adequada é:
\begin{enumerate}[label=\Alph*)]
\item publicar porque LAI sempre prevalece sobre LGPD;
\item impedir qualquer publicação porque LGPD sempre prevalece;
\item avaliar finalidade, base legal, necessidade, restrições, minimização e hipóteses de acesso;
\item publicar somente se a base estiver em CSV;
\item anonimizar apenas o nome, independentemente de outros identificadores.
\end{enumerate}
\end{questao}

\begin{questao}{\nivel{N2} 14.25 — Cebraspe/Certo ou Errado}
Se duas áreas usam definições diferentes de ``contrato ativo'', um glossário de negócio com owner e processo de aprovação pode reduzir inconsistência semântica.
\end{questao}

\subsection*{N3 — Integração de concurso}

\begin{questao}{\nivel{N3} 14.26 — FGV/assertivas COBIT}
Considere:

I. EDM é governança.

II. APO, BAI, DSS e MEA são gestão.

III. Uma organização deve implementar todos os objetivos COBIT com a mesma prioridade para preservar aderência ao framework.

Assinale a alternativa correta.
\begin{enumerate}[label=\Alph*)]
\item Apenas I.
\item Apenas II.
\item Apenas I e II.
\item Apenas II e III.
\item I, II e III.
\end{enumerate}
\end{questao}

\begin{questao}{\nivel{N3} 14.27 — Cebraspe/Certo ou Errado}
Uma atividade de monitoramento pode pertencer à gestão ou à governança; o enquadramento depende do objeto, do nível decisório e da responsabilidade envolvida.
\end{questao}

\begin{questao}{\nivel{N3} 14.28 — FCC/ITIL}
Uma organização possui alto volume de incidentes recorrentes. A equipe mede apenas tempo de restauração e não registra problemas conhecidos. O efeito provável é:
\begin{enumerate}[label=\Alph*)]
\item boa maturidade de problem management;
\item restauração rápida sem redução de recorrência;
\item eliminação de causas;
\item melhoria garantida de disponibilidade;
\item substituição de incidentes por requisições.
\end{enumerate}
\end{questao}

\begin{questao}{\nivel{N3} 14.29 — FGV/value stream}
Um fluxo de valor possui várias aprovações manuais redundantes, 70\% do lead time em espera e baixa taxa de rejeição. Antes de automatizar as aprovações, a melhor ação é:
\begin{enumerate}[label=\Alph*)]
\item automatizar todas sem questionar;
\item avaliar valor e risco de cada etapa, eliminando controles redundantes antes de automatizar;
\item aumentar o número de aprovações;
\item proibir automação;
\item migrar o processo para Scrum.
\end{enumerate}
\end{questao}

\begin{questao}{\nivel{N3} 14.30 — Cebraspe/Certo ou Errado}
Um indicador de disponibilidade cuja regra de exclusão de indisponibilidades pode ser alterada manualmente após o período de medição apresenta fragilidade de governança, mesmo que a fórmula matemática seja correta.
\end{questao}

\begin{questao}{\nivel{N3} 14.31 — FCC/portfolio}
Duas iniciativas competem por orçamento. A primeira é obrigação regulatória de baixo benefício financeiro; a segunda tem alto ROI, mas não é obrigatória. Um modelo de priorização maduro deve:
\begin{enumerate}[label=\Alph*)]
\item ordenar somente por ROI;
\item considerar obrigação legal, risco, valor, dependências e capacidade;
\item escolher sempre a mais barata;
\item excluir critérios qualitativos;
\item usar apenas prazo.
\end{enumerate}
\end{questao}

\begin{questao}{\nivel{N3} 14.32 — FGV/PMBOK}
Um projeto foi entregue no prazo e orçamento, mas apenas 20\% dos usuários adotaram a solução e os benefícios esperados não se materializaram. A melhor conclusão é:
\begin{enumerate}[label=\Alph*)]
\item o projeto foi plenamente bem-sucedido porque cumpriu prazo e custo;
\item a análise de sucesso deve incluir valor, benefícios e adoção além da entrega;
\item adoção é tema exclusivamente operacional;
\item benefícios não devem ser medidos;
\item a falha prova que métodos ágeis são inadequados.
\end{enumerate}
\end{questao}

\begin{questao}{\nivel{N3} 14.33 — Cebraspe/Certo ou Errado]
A existência de 40 KPIs sem definição de fonte, fórmula ou owner pode reduzir, e não aumentar, a qualidade da governança de desempenho.
\end{questao}

\begin{questao}{\nivel{N3} 14.34 — FCC/BPMN}
Um gateway paralelo diverge em dois caminhos. Para sincronizar ambos antes de prosseguir, utiliza-se, em regra:
\begin{enumerate}[label=\Alph*)]
\item outro gateway paralelo convergente;
\item gateway exclusivo;
\item message flow;
\item evento de início;
\item pool adicional.
\end{enumerate}
\end{questao}

\begin{questao}{\nivel{N3} 14.35 — FGV/compliance}
Uma entidade atende aos requisitos mínimos de uma norma, mas possui risco residual acima do apetite institucional. A decisão correta é:
\begin{enumerate}[label=\Alph*)]
\item conformidade encerra a análise;
\item controles adicionais ou tratamento de risco podem ser necessários apesar da conformidade;
\item o apetite deve ser aumentado automaticamente;
\item certificação elimina risco residual;
\item risco e compliance são sinônimos.
\end{enumerate}
\end{questao}

\begin{questao}{\nivel{N3} 14.36 — Cebraspe/Certo ou Errado}
Um fornecedor certificado pode reduzir parte do risco de sourcing, mas o cliente continua responsável por suas próprias configurações, acessos, dados e obrigações não cobertas pelo escopo do fornecedor.
\end{questao}

\begin{questao}{\nivel{N3} 14.37 — FCC/DAMA}
Uma organização possui múltiplas versões de ``cliente'', cada sistema com identificadores e regras diferentes. A disciplina mais diretamente envolvida em resolver a visão corporativa da entidade é:
\begin{enumerate}[label=\Alph*)]
\item master/reference data management, sob governança de dados;
\item incident management;
\item gestão de configuração de rede;
\item BPMN exclusivamente;
\item RAID.
\end{enumerate}
\end{questao}

\begin{questao}{\nivel{N3} 14.38 — FGV/risco}
A área de TI classifica um risco como aceitável porque o custo do controle é alto, mas não existe critério de apetite ou autorização da governança. A fragilidade é:
\begin{enumerate}[label=\Alph*)]
\item aceitação sem autoridade e critério formal;
\item excesso de controle;
\item ausência de ITIL;
\item falta de indicador financeiro apenas;
\item uso de análise qualitativa.
\end{enumerate}
\end{questao}

\begin{questao}{\nivel{N3} 14.39 — Cebraspe/Certo ou Errado}
A adoção de Scrum em uma equipe não elimina a necessidade de governança de portfólio, orçamento, risco e arquitetura em nível organizacional.
\end{questao}

\begin{questao}{\nivel{N3} 14.40 — FCC/BSC}
Um mapa estratégico apresenta dezenas de indicadores sem relação explícita com objetivos. A principal correção é:
\begin{enumerate}[label=\Alph*)]
\item aumentar ainda mais os indicadores;
\item ligar indicadores a objetivos e relações de causa/resultado relevantes;
\item remover metas;
\item usar somente indicadores financeiros;
\item substituir BSC por service desk.
\end{enumerate}
\end{questao}

\subsection*{N4 — Auditoria / avançado}

\begin{questao}{\nivel{N4} 14.41 — Cebraspe/caso integrado}
O PDTI possui 80 iniciativas, mas nenhuma está vinculada a objetivos institucionais, responsáveis ou indicadores. Julgue o item:

Ainda que todas as iniciativas sejam executadas no prazo, a organização pode não demonstrar alinhamento estratégico nem geração de valor.
\end{questao}

\begin{questao}{\nivel{N4} 14.42 — FGV/KPI}
A diretoria de TI recebe bônus quando 95\% dos projetos terminam no prazo. Para atingir a meta, equipes reduzem escopo perto do encerramento e classificam entregas parciais como concluídas. A melhor resposta de governança é:
\begin{enumerate}[label=\Alph*)]
\item elevar a meta para 99\%;
\item balancear prazo com escopo entregue, benefícios, qualidade, adoção e transparência das mudanças;
\item eliminar qualquer indicador;
\item manter a meta e punir apenas atrasos;
\item medir somente custo.
\end{enumerate}
\end{questao}

\begin{questao}{\nivel{N4} 14.43 — FCC/COBIT}
Uma organização implementa todos os objetivos COBIT igualmente, embora tenha poucos recursos e riscos concentrados em cibersegurança e dados. O principal problema é:
\begin{enumerate}[label=\Alph*)]
\item excesso de alinhamento;
\item ausência de tailoring e priorização pelos fatores de desenho;
\item uso de COBIT em órgão público;
\item falta de ITIL;
\item baixa quantidade de processos.
\end{enumerate}
\end{questao}

\begin{questao}{\nivel{N4} 14.44 — Cebraspe/Certo ou Errado}
Um comitê que aprova investimentos de TI, mas não acompanha posteriormente benefícios, riscos ou desempenho, executa apenas parte do ciclo de governança.
\end{questao}

\begin{questao}{\nivel{N4} 14.45 — FGV/BPMN e auditoria}
Um processo de concessão de acesso possui cinco aprovações sequenciais idênticas, mas nenhuma verifica atributo diferente. O lead time é alto e quase nenhuma solicitação é rejeitada. A recomendação mais adequada é:
\begin{enumerate}[label=\Alph*)]
\item eliminar todas as aprovações sem análise;
\item avaliar risco e propósito de cada aprovação, consolidando controles redundantes e preservando accountability;
\item automatizar as cinco etapas sem revisão;
\item adicionar uma sexta aprovação;
\item substituir o processo por e-mail.
\end{enumerate}
\end{questao}

\begin{questao}{\nivel{N4} 14.46 — Cebraspe/Certo ou Errado}
Uma meta de disponibilidade pode ser formalmente atingida e, ainda assim, ser pouco útil à governança se a definição do indicador excluir períodos críticos para o negócio.
\end{questao}

\begin{questao}{\nivel{N4} 14.47 — FCC/governança de dados}
Dois dashboards oficiais apresentam números distintos porque cada área aplica regra própria de ``despesa liquidada''. O controle prioritário é:
\begin{enumerate}[label=\Alph*)]
\item criar um terceiro dashboard;
\item estabelecer definição de negócio oficial, owner, lineage e reconciliação;
\item aumentar a capacidade do banco;
\item eliminar metadados;
\item usar apenas planilhas.
\end{enumerate}
\end{questao}

\begin{questao}{\nivel{N4} 14.48 — FGV/risco e alta administração}
A área técnica decide unilateralmente aceitar risco de indisponibilidade de serviço crítico porque a mitigação custa caro. A principal falha é:
\begin{enumerate}[label=\Alph*)]
\item inexistência de autorização de governança compatível com o apetite e impacto institucional;
\item uso de análise de custo;
\item participação da TI;
\item inexistência de Scrum;
\item ausência de SLA externo.
\end{enumerate}
\end{questao}

\begin{questao}{\nivel{N4} 14.49 — Cebraspe/Certo ou Errado}
A transparência da LAI e a proteção da LGPD não devem ser tratadas como comandos mutuamente excludentes; a decisão de divulgação exige análise do dado, finalidade e regime jurídico aplicável.
\end{questao}

\begin{questao}{\nivel{N4} 14.50 — FGV/matriz de achado}
A auditoria identifica 42 KPIs de TI; 19 não possuem fonte definida, 11 são calculados manualmente e 8 mudaram de fórmula sem registro. Qual achado é mais bem formulado?
\begin{enumerate}[label=\Alph*)]
\item ``Há KPIs demais'';
\item ``Os indicadores são inúteis'';
\item ``A ausência de governança de definição, fonte e mudança dos KPIs reduz confiabilidade e comparabilidade da medição; recomenda-se catálogo de indicadores, owners, versionamento de fórmula e automação/reconciliação'';
\item ``Deve-se usar apenas BSC'';
\item ``Basta reduzir todos os KPIs pela metade''.
\end{enumerate}
\end{questao}
